\def\uadectitulo{Fisiología humana}
\documentclass[12pt,letterpaper,oneside]{article}
\usepackage[T1]{fontenc}
\usepackage[utf8]{inputenc}
\usepackage[spanish,es-nodecimaldot]{babel}
\usepackage{lmodern}
\usepackage[margin=2.5cm]{geometry}
\usepackage{amsmath,amssymb,graphicx,booktabs,longtable,array,enumitem}
\usepackage[most]{tcolorbox}
\usepackage{tikz}
\usetikzlibrary{babel,arrows.meta,positioning}
\usepackage{setspace}
\usepackage[hidelinks]{hyperref}
\onehalfspacing
\setlength{\emergencystretch}{3em}
\providecommand{\uadectitulo}{Material de estudio}
\providecommand{\uadecautor}{Material de estudio AulaTeX}
\newcommand{\portadaUAdeC}{%
  \hypersetup{pageanchor=false}
  \begin{titlepage}
  \centering
  {\Large Universidad Autónoma de Coahuila\par}
  \vspace{0.7cm}
  {\large Preparación para ingreso a Médico Cirujano\par}
  \vspace{1.8cm}
  {\LARGE\bfseries\uadectitulo\par}
  \vspace{1cm}
  {\large Actividad 1: síntesis y representación conceptual\par}
  {\large Actividad 2: práctica y soluciones justificadas\par}
  \vfill
  {\uadecautor\par}
  {Revisión: 10 de octubre de 2026\par}
  \vspace{0.8cm}
  \begin{minipage}{0.87\linewidth}
  \small Material formativo local. No constituye una guía, plantilla o examen oficial de UAdeC o CENEVAL. La cobertura y dificultad no están certificadas frente a una convocatoria de ingreso. No sustituye orientación ni atención clínica.
  \end{minipage}
  \end{titlepage}
  \hypersetup{pageanchor=true}
  \tableofcontents\clearpage
  \section*{Alcance y procedencia}
  Las actividades siguen las consignas locales conservadas en esta materia. Los reactivos son de práctica y las respuestas son propuestas razonadas, no resultados de un intento del aspirante. Las fuentes de apoyo se identifican al final. Los PDF bibliográficos locales se inventarían por archivo y hash; su disponibilidad no acredita lectura íntegra ni permiso de redistribución.
}
\begin{document}
\portadaUAdeC
\section{Actividad 1: síntesis de investigación}
\subsection{Homeostasis}
Homeostasis es regulación dinámica dentro de rangos compatibles con función, no constancia absoluta. La retroalimentación negativa contrarresta una desviación y la positiva amplifica un proceso que necesita un mecanismo de término, como el parto. Sensor, centro integrador y efector son roles funcionales que pueden estar combinados en un mismo órgano.
\subsection{Circulación y respiración}
Sístole es contracción y diástole relajación; las cámaras no se contraen todas al mismo tiempo. El nodo sinoauricular inicia normalmente la actividad cardíaca. La presión depende de gasto cardíaco y resistencia, entre otros factores. Ventilación mueve aire; difusión intercambia gases a favor de gradientes de presión parcial y perfusión transporta sangre. La mayor parte del CO2 circula como bicarbonato, aunque existen otras formas de transporte.
\subsection{Endocrino y renal}
Hipotálamo e hipófisis coordinan varios ejes, pero no todas las glándulas dependen de un único eje. TSH estimula tiroides; insulina favorece reducción de glucemia y glucagón movilización de reservas. La filtración renal se inicia en el corpúsculo renal; túbulos realizan reabsorción y secreción. Excreción resulta de filtración menos reabsorción más secreción. La nefrona es la unidad funcional y no debe confundirse con un uréter ni con toda la pelvis renal.
\subsection{Representación conceptual y cuadro de síntesis}
\begin{longtable}{p{0.27\linewidth}p{0.65\linewidth}}
\toprule Concepto & Relación o función\\\midrule
Variable regulada & Temperatura o glucemia\\
Sensor & Detecta desviación\\
Integración & Compara y coordina\\
Efector & Modifica el proceso\\
Respuesta negativa & Reduce la desviación inicial\\
\bottomrule
\end{longtable}
\begin{figure}[htbp]
\centering
\begin{tikzpicture}[node distance=5mm, every node/.style={draw,rounded corners=2pt,text width=10cm,align=center,font=\small}]
\node (step0) {Variable regulada: Temperatura o glucemia};
\node[below=of step0] (step1) {Sensor: Detecta desviación};
\draw[-{Stealth}] (step0) -- (step1);
\node[below=of step1] (step2) {Integración: Compara y coordina};
\draw[-{Stealth}] (step1) -- (step2);
\node[below=of step2] (step3) {Efector: Modifica el proceso};
\draw[-{Stealth}] (step2) -- (step3);
\node[below=of step3] (step4) {Respuesta negativa: Reduce la desviación inicial};
\draw[-{Stealth}] (step3) -- (step4);
\end{tikzpicture}
\caption{Esquema didáctico original; las flechas indican orden de lectura o relación conceptual, no causalidad universal.}
\end{figure}

\clearpage



\section{Módulo Teórico: Fisiología de Alto Rendimiento}
\subsection{Homeostasis y Medio Interno}
La \textbf{homeostasis} es el mantenimiento de condiciones casi constantes en el medio interno. Se logra mediante sistemas de retroalimentación:
\begin{itemize}
    \item \textbf{Retroalimentación Negativa:} El producto final inhibe o frena el estímulo inicial. Es el mecanismo más común. Ejemplos: regulación de la temperatura, glucemia por insulina, presión arterial.
    \item \textbf{Retroalimentación Positiva:} El producto final amplifica el estímulo inicial, llevando el proceso a su culminación. Ejemplos: coagulación sanguínea, contracciones uterinas (oxitocina), potencial de acción (apertura de canales de sodio).
\end{itemize}

\subsection{Fisiología Cardiovascular}
\begin{itemize}
    \item \textbf{Sistema de Conducción:} Nodo Sinoauricular (SA) [marcapasos natural, 60-100 lpm] $\rightarrow$ Vías internodales $\rightarrow$ Nodo Auriculoventricular (AV) [retraso de 0.1s para permitir el llenado ventricular] $\rightarrow$ Haz de His $\rightarrow$ Fibras de Purkinje.
    \item \textbf{Ciclo Cardíaco:}
    \begin{itemize}
        \item \textit{Sístole (Contracción):} Cierre de válvulas AV (Mitral y Tricúspide) = \textbf{Primer ruido cardíaco (R1)}. Eyección de sangre.
        \item \textit{Diástole (Relajación):} Cierre de válvulas Semilunares (Aórtica y Pulmonar) = \textbf{Segundo ruido cardíaco (R2)}. Llenado ventricular.
    \end{itemize}
\end{itemize}
\begin{tcolorbox}[colback=blue!5!white,colframe=blue!75!black,title=Nota de estudio]
El \textbf{gasto cardíaco} es el volumen de sangre bombeado por minuto (aprox. 5 L/min). Se calcula multiplicando la \textbf{Frecuencia Cardíaca (FC)} por el \textbf{Volumen Sistólico (VS)}.
\end{tcolorbox}

\subsection{Fisiología Respiratoria}
\begin{itemize}
    \item \textbf{Ventilación:} Inspiración (activa, contrae el diafragma) y Espiración (pasiva, relajación elástica).
    \item \textbf{Volúmenes:} \textit{Volumen Corriente (Tidal)} = 500 mL (aire respirado en reposo).
    \item \textbf{Intercambio gaseoso:} Por difusión simple a favor de un gradiente de presiones parciales. El factor surfactante (producido por neumocitos tipo II) reduce la tensión superficial alveolar y evita el colapso.
\end{itemize}

\subsection{Sistema Endocrino}
Hormonas principales y funciones:
\begin{itemize}
    \item \textbf{Hipotálamo:} Controla a la hipófisis mediante factores liberadores.
    \item \textbf{Neurohipófisis (Lóbulo posterior):} Almacena \textit{Oxitocina} y \textit{Vasopresina (ADH)}.
    \item \textbf{Adenohipófisis (Lóbulo anterior):} Secreta TSH, ACTH, FSH, LH, GH, Prolactina.
    \item \textbf{Páncreas:} Islotes de Langerhans. Células alfa (Glucagón = hiperglucemiante), Células beta (Insulina = hipoglucemiante).
\end{itemize}

\subsection{Fisiología Renal}
La nefrona es la unidad funcional (aprox. 1 millón por riñón).
\begin{itemize}
    \item \textbf{Glomérulo:} Filtración de sangre. No se filtran proteínas grandes ni eritrocitos.
    \item \textbf{Túbulo Contorneado Proximal (TCP):} Reabsorbe el 65\% del agua, Na+ y el 100\% de glucosa y aminoácidos sanos.
    \item \textbf{Asa de Henle:} Concentra la orina.
    \item \textbf{Túbulo Contorneado Distal (TCD) y Colector:} Regulados por hormonas (ADH reabsorbe agua; Aldosterona reabsorbe Na+ y excreta K+).
\end{itemize}

\newpage
\section{Banco de Reactivos: Actividad 2: banco original de práctica}
\begin{enumerate}[label=\textbf{\arabic*.}]
    \item \textbf{Un paciente presenta hipotensión. En respuesta, su cuerpo inicia una serie de procesos para elevar la presión arterial hasta volver a niveles normales, momento en el cual los mecanismos se detienen. Este es un ejemplo de:}
    \begin{enumerate}
        \item[A)] Retroalimentación positiva
        \item[B)] Retroalimentación negativa
        \item[C)] Alostasis
        \item[D)] Feed-forward
    \end{enumerate}
    \textit{Respuesta: B. El estímulo inicial (hipotensión) es revertido por el sistema, el cual se apaga al lograr el objetivo. Es el clásico feedback negativo homeostático.}

    \item \textbf{¿Cuál de las siguientes estructuras es considerada el marcapasos fisiológico del corazón?}
    \begin{enumerate}
        \item[A)] Nodo auriculoventricular (AV)
        \item[B)] Fibras de Purkinje
        \item[C)] Haz de His
        \item[D)] Nodo sinoauricular (SA)
    \end{enumerate}
    \textit{Respuesta: D. El nodo SA en la aurícula derecha tiene la despolarización intrínseca más rápida, dictando el ritmo de contracción de todo el corazón.}

    \item \textbf{Durante la auscultación cardíaca de un paciente sano, el médico escucha el primer ruido cardíaco (R1). ¿Qué evento mecánico produce este ruido?}
    \begin{enumerate}
        \item[A)] Cierre de las válvulas semilunares (aórtica y pulmonar)
        \item[B)] Apertura de las válvulas auriculoventriculares
        \item[C)] Cierre de las válvulas auriculoventriculares (mitral y tricúspide)
        \item[D)] Llenado ventricular rápido
    \end{enumerate}
    \textit{Respuesta: C. El R1 ocurre al inicio de la sístole ventricular por el cierre de las válvulas AV. El R2 ocurre por el cierre de las semilunares.}

    \item \textbf{En el electrocardiograma (ECG), ¿qué representa el complejo QRS?}
    \begin{enumerate}
        \item[A)] Despolarización auricular
        \item[B)] Despolarización ventricular
        \item[C)] Repolarización ventricular
        \item[D)] Repolarización auricular
    \end{enumerate}
    \textit{Respuesta: B. La onda P es despolarización auricular. El QRS es despolarización ventricular. La onda T es repolarización ventricular.}

    \item \textbf{¿Qué célula pulmonar es responsable de producir el factor surfactante, vital para evitar el colapso alveolar durante la espiración?}
    \begin{enumerate}
        \item[A)] Neumocito tipo I
        \item[B)] Macrófago alveolar
        \item[C)] Célula caliciforme
        \item[D)] Neumocito tipo II
    \end{enumerate}
    \textit{Respuesta: D. El neumocito tipo II secreta surfactante. El tipo I forma la barrera de intercambio gaseoso. Los macrófagos fagocitan partículas.}

    \item \textbf{Un volumen respiratorio normal en reposo, en cada ciclo de inspiración y espiración, es de aproximadamente 500 mL. ¿Cómo se denomina este volumen?}
    \begin{enumerate}
        \item[A)] Volumen de reserva inspiratoria
        \item[B)] Volumen residual
        \item[C)] Volumen corriente (Tidal)
        \item[D)] Capacidad vital
    \end{enumerate}
    \textit{Respuesta: C. El volumen corriente (Tidal) es la cantidad de aire que entra o sale en una respiración tranquila.}

    \item \textbf{¿En qué región de la nefrona se reabsorbe normalmente casi toda la glucosa filtrada, mientras no se supere la capacidad de transporte tubular?}
    \begin{enumerate}
        \item[A)] Glomérulo
        \item[B)] Túbulo contorneado proximal
        \item[C)] Asa de Henle
        \item[D)] Túbulo colector
    \end{enumerate}
    \textit{Respuesta: B. El TCP es la principal zona de reabsorción obligatoria masiva de nutrientes (glucosa, aminoácidos) e iones.}

    \item \textbf{Ante un estímulo apropiado de osmolaridad o volumen, la neurohipófisis libera una hormona sintetizada en el hipotálamo que aumenta la reabsorción de agua en los conductos colectores. ¿Cuál es?}
    \begin{enumerate}
        \item[A)] Aldosterona
        \item[B)] Oxitocina
        \item[C)] Hormona Antidiurética (ADH / Vasopresina)
        \item[D)] Péptido Natriurético Auricular (ANP)
    \end{enumerate}
    \textit{Respuesta: C. La ADH inserta acuaporinas en el túbulo colector para retener agua. La aldosterona retiene Na+ y secreta K+. El ANP promueve la pérdida de sodio y agua.}

    \item \textbf{Las hormonas esteroideas (como el cortisol o la testosterona) difieren de las hormonas peptídicas (como la insulina) principalmente porque:}
    \begin{enumerate}
        \item[A)] Actúan exclusivamente mediante segundos mensajeros en la membrana.
        \item[B)] Pueden atravesar la membrana plasmática y unirse a receptores intracelulares.
        \item[C)] Viajan disueltas libremente en el plasma sin requerir proteínas transportadoras.
        \item[D)] Son secretadas únicamente por la hipófisis.
    \end{enumerate}
    \textit{Respuesta: B. Por ser derivadas del colesterol (lipofílicas), las hormonas esteroideas atraviesan la membrana lipídica celular y actúan directo en el núcleo.}

    \item \textbf{Al medir la presión arterial, la lectura "120/80 mmHg" indica:}
    \begin{enumerate}
        \item[A)] Presión diastólica de 120 y sistólica de 80.
        \item[B)] Presión sistólica de 120 y diastólica de 80.
        \item[C)] Frecuencia cardíaca de 120 y gasto cardíaco de 80.
        \item[D)] Presión venosa central de 120 y presión arterial de 80.
    \end{enumerate}
    \textit{Respuesta: B. El valor mayor (120) es la presión máxima o sistólica; el valor menor (80) es la presión mínima o diastólica de relajación.}

    \item \textbf{La enzima pepsina, fundamental para la digestión de proteínas, se secreta en el estómago en forma de pepsinógeno inactivo. ¿Qué estimula su conversión a la forma activa?}
    \begin{enumerate}
        \item[A)] El bicarbonato pancreático
        \item[B)] La bilis
        \item[C)] El Ácido Clorhídrico (HCl)
        \item[D)] La lipasa gástrica
    \end{enumerate}
    \textit{Respuesta: C. El medio ácido (HCl, secretado por las células parietales) escinde el pepsinógeno para formar pepsina activa.}

    \item \textbf{¿Qué hormona es secretada por las células alfa de los islotes de Langerhans pancreáticos durante el ayuno para elevar la glucemia?}
    \begin{enumerate}
        \item[A)] Insulina
        \item[B)] Glucagón
        \item[C)] Somatostatina
        \item[D)] Adrenalina
    \end{enumerate}
    \textit{Respuesta: B. El glucagón es hiperglucemiante (activa glucogenólisis en el hígado). Las células beta producen insulina (hipoglucemiante).}

    \item \textbf{¿Cuál es la función principal de la bilis en el proceso digestivo?}
    \begin{enumerate}
        \item[A)] Digerir proteínas a aminoácidos.
        \item[B)] Emulsificar las grasas para facilitar su digestión por las lipasas.
        \item[C)] Acidificar el quimo al llegar al duodeno.
        \item[D)] Absorber la vitamina B12.
    \end{enumerate}
    \textit{Respuesta: B. Las sales biliares favorecen la emulsificación de gotas lipídicas y aumentan la superficie accesible a las lipasas. También participan en micelas mixtas que facilitan la absorción; emulsión y micela no son sinónimos. La bilis no es una proteasa ni acidifica el quimo.}

    \item \textbf{En fisiología neuromuscular, ¿cuál es el neurotransmisor liberado en la placa motora terminal para inducir la contracción del músculo esquelético?}
    \begin{enumerate}
        \item[A)] Adrenalina
        \item[B)] Serotonina
        \item[C)] Acetilcolina
        \item[D)] Dopamina
    \end{enumerate}
    \textit{Respuesta: C. La acetilcolina (ACh) es el neurotransmisor exclusivo de la unión neuromuscular somática (músculo esquelético).}

    \item \textbf{El eritrocito (glóbulo rojo) es responsable de transportar oxígeno debido a que contiene:}
    \begin{enumerate}
        \item[A)] Mioglobina
        \item[B)] Hemoglobina
        \item[C)] Inmunoglobulina
        \item[D)] Albúmina
    \end{enumerate}
    \textit{Respuesta: B. La hemoglobina es la proteína férrica intracelular del eritrocito que fija el O2. La mioglobina está en el músculo.}

    \item \textbf{¿Cuál es el efecto del sistema nervioso simpático sobre el corazón?}
    \begin{enumerate}
        \item[A)] Disminuir la frecuencia cardíaca
        \item[B)] Reducir la fuerza de contracción
        \item[C)] Aumentar la frecuencia cardíaca y la fuerza de contracción
        \item[D)] Vasodilatar las arterias coronarias induciendo bradicardia
    \end{enumerate}
    \textit{Respuesta: C. El simpático (lucha y huida, adrenalina/noradrenalina) causa taquicardia (cronotropismo +) e inotropismo +.}

    \item \textbf{El intercambio de gases entre los capilares sistémicos y las células de los tejidos del cuerpo (donde el O2 entra a la célula y el CO2 pasa a la sangre) se conoce como:}
    \begin{enumerate}
        \item[A)] Ventilación pulmonar
        \item[B)] Respiración celular
        \item[C)] Respiración externa
        \item[D)] Respiración interna (tisular)
    \end{enumerate}
    \textit{Respuesta: D. La respiración externa ocurre en los alvéolos pulmonares. La interna u ocurre a nivel de todos los tejidos periféricos.}

    \item \textbf{¿Qué órgano endocrino se encarga de regular la tasa de metabolismo basal mediante la secreción de las hormonas T3 y T4?}
    \begin{enumerate}
        \item[A)] Hipotálamo
        \item[B)] Glándula Tiroides
        \item[C)] Glándula Paratiroides
        \item[D)] Médula suprarrenal
    \end{enumerate}
    \textit{Respuesta: B. La tiroides secreta tiroxina (T4) y triyodotironina (T3), responsables de controlar el índice metabólico celular.}

    \item \textbf{El reflejo miotático, responsable del "tirón" de la rodilla al percutir el tendón rotuliano, es un ejemplo de un reflejo medular porque:}
    \begin{enumerate}
        \item[A)] Su integración ocurre en la corteza motora cerebral.
        \item[B)] Su centro de integración se localiza en la médula espinal, sin necesidad de comando encefálico inmediato.
        \item[C)] Depende de los nervios craneales.
        \item[D)] Es procesado en el cerebelo.
    \end{enumerate}
    \textit{Respuesta: B. Los reflejos simples y rápidos se integran a nivel medular para una respuesta defensiva veloz, aunque luego informan al cerebro.}

    \item \textbf{Durante una hemorragia severa (hipovolemia), se activa el sistema Renina-Angiotensina-Aldosterona (SRAA). El efecto final de este sistema es:}
    \begin{enumerate}
        \item[A)] Vasodilatación generalizada
        \item[B)] Retención de sodio y agua, con aumento de la presión arterial
        \item[C)] Excreción masiva de orina (diuresis)
        \item[D)] Disminución de la frecuencia cardíaca
    \end{enumerate}
    \textit{Respuesta: B. El SRAA busca restaurar el volumen sanguíneo y la presión (Angiotensina II es un potente vasoconstrictor y estimula a la aldosterona para retener sodio y agua).}

\end{enumerate}

\section{Conclusiones}
El material permite repasar conceptos y contrastar respuestas. La clave no acredita desempeño del aspirante ni aprobación institucional. Deben revisarse la convocatoria y el temario aplicables antes de usarlo como plan definitivo de ingreso.
\clearpage\section{Fuentes de apoyo y alcance de consulta}
OpenStax. Anatomy and Physiology 2e: 1.5 Homeostasis. Consulta: 10 de octubre de 2026. \url{https://openstax.org/books/anatomy-and-physiology-2e/pages/1-5-homeostasis}\par Sección consultada: regulación y tipos de retroalimentación; otros sistemas requieren cotejo específico.\par\medskip
Los originales bibliográficos disponibles se identifican por nombre de archivo y SHA256 en el inventario local. No se atribuyen ediciones, páginas ni lecturas no comprobadas.
\end{document}
